\title{Chain-of-Thought Prompting Elicits Reasoning in Large Language Models}

\begin{abstract}
We investigate how producing a \textit{chain of thought}—a sequence of intermediate reasoning steps—greatly enhances the capability of large language models to handle complex reasoning. Specifically, we demonstrate that such reasoning skills arise naturally in sufficiently large language models through a straightforward approach called \textit{chain-of-thought prompting}, which involves presenting a few chain of thought examples as prompts.

Tests conducted on three large language models reveal that chain-of-thought prompting boosts performance across a variety of arithmetic, commonsense, and symbolic reasoning tasks. The improvements observed can be remarkable. For example, prompting a PaLM 540B model with only eight chain-of-thought examples attains state-of-the-art results on the GSM8K dataset of math word problems, outperforming even fine-tuned GPT-3 models that use a verifier.
\end{abstract}

\section{Introduction}
The field of NLP has recently undergone significant transformation due to language models \citep[][\textit{inter alia}]{peters-etal-2018-deep,devlin-etal-2019-bert,brown2020language}. Increasing the scale of language models has been shown to yield various advantages, including enhanced performance and sample efficiency \citep[\textit{inter alia}]{kaplan2020scaling,brown2020language}. Nonetheless, simply enlarging model size has not been enough to achieve strong results on difficult tasks such as arithmetic, commonsense, and symbolic reasoning \citep{rae2021scaling}.

In this work, we investigate how to unlock the reasoning potential of large language models through a simple technique inspired by two insights. First, methods for arithmetic reasoning can benefit from generating natural language explanations that lead to the final answer. Previous research has enabled models to produce natural language intermediate steps by training from scratch \citep{ling-etal-2017-program} or by fine-tuning pretrained models \citep{cobbe2021training}, alongside neuro-symbolic approaches that employ formal languages instead of natural language \citep{roy-roth-2015-solving, chiang-chen-2019-semantically,amini-etal-2019-mathqa, chen2019neural}. Second, large language models present an exciting opportunity for in-context few-shot learning through \emph{prompting}. Instead of fine-tuning a distinct language model checkpoint for every new task, one can simply ``prompt'' the model with a few input-output pairs illustrating the task. Impressively, this method has been effective across various straightforward question-answering tasks \citep{brown2020language}.

Nonetheless, both of these concepts have significant drawbacks. In the case of rationale-augmented training and finetuning techniques, generating an extensive collection of high-quality rationales is resource-intensive and considerably more complex than the straightforward input--output pairs typically employed in standard machine learning. On the other hand, the conventional few-shot prompting approach utilized in \citet{brown2020language} tends to perform inadequately on tasks that demand reasoning capabilities and often shows limited improvement even as the size of the language model increases \citep{rae2021scaling}. In this study, we integrate the advantages of these two methods while circumventing their respective shortcomings. Specifically, we investigate the capability of language models to execute few-shot prompting for reasoning tasks using prompts composed of triples: $\langle$input, \emph{chain of thought}, output$\rangle$.  
A \emph{chain of thought} consists of a sequence of intermediate natural language reasoning steps that culminate in the final answer, and we designate this technique as \textit{chain-of-thought prompting}. An example of such a prompt is provided in \cref{fig:pull-figure}.  

We conduct empirical tests across benchmarks involving arithmetic, commonsense, and symbolic reasoning, demonstrating that chain-of-thought prompting surpasses standard prompting, at times by a substantial margin. \cref{fig:pull-bar-chart} depicts one such outcome—in the GSM8K benchmark of math word problems \citep{cobbe2021training}, chain-of-thought prompting with \palm{} 540B significantly outperforms conventional prompting and establishes a new state-of-the-art result. Employing a prompting-only approach is valuable because it does not necessitate a large labeled training set and allows a single model checkpoint to tackle numerous tasks without compromising generality. This research highlights how large language models can learn from just a few examples that include natural language descriptions of the task (as opposed to solely relying on extensive training datasets to implicitly learn input-output patterns).

\section{Chain-of-Thought Prompting} When tackling a complex reasoning problem, such as a multi-step math word problem, individuals often reflect on their own thought process. It is common to break down the problem into smaller intermediate steps and address each one before arriving at the final solution: \textit{``After Jane gives 2 flowers to her mom she has 10 $\ldots$ then after she gives 3 to her dad she will have 7 $\ldots$ so the answer is 7.''} The objective of this paper is to equip language models with the capability to produce a similar \textit{chain of thought}—a logically connected sequence of intermediate reasoning steps that culminate in the final answer to a problem. We demonstrate that language models of adequate size can generate these chains of thought when chain-of-thought reasoning examples are included in the few-shot prompting exemplars.

\cref{fig:pull-figure} presents an instance where a model generates a chain of thought to solve a math word problem it would otherwise have answered incorrectly. The chain of thought here resembles a solution and can be understood as such, but we choose to term it a chain of thought to emphasize that it imitates a stepwise reasoning process leading to the answer (and also, solutions or explanations typically follow \textit{after} the final answer \citep[][\textit{inter alia}]{narang2020wt5,wiegreffe2021reframing,lampinen2022can}).

Chain-of-thought prompting offers several appealing advantages as a method to enhance reasoning capabilities in language models.

\begin{enumerate}[topsep=1pt,itemsep=0ex]
    \item To begin with, chain of thought inherently enables models to break down complex multi-step problems into smaller intermediate steps, allowing for additional computational resources to be devoted to problems that necessitate more reasoning stages.
    \item Next, a chain of thought offers a transparent view into the model’s reasoning process, indicating how it may have reached a specific conclusion and offering opportunities to identify and correct errors in the reasoning trajectory (although fully capturing the computations a model uses to support an answer remains an unresolved challenge).
    \item Moreover, chain-of-thought reasoning can be applied to tasks such as solving math word problems, commonsense reasoning, and symbolic manipulation, and is theoretically applicable to any task that humans can solve using language.
    \item Lastly, chain-of-thought reasoning can be easily induced in sufficiently large pretrained language models simply by including chain-of-thought sequences within the few-shot prompting examples.
\end{enumerate}

In our empirical studies, we demonstrate the effectiveness of chain-of-thought prompting on arithmetic reasoning (\cref{sec:arithmetic-reasoning}), commonsense reasoning (\cref{sec:commonsense-reasoning}), and symbolic reasoning (\cref{sec:symbolic-reasoning}).

\section{Arithmetic Reasoning}\label{sec:arithmetic-reasoning}
We start by examining math word problems of the type illustrated in \cref{fig:pull-figure}, which evaluate the arithmetic reasoning capabilities of language models. Although straightforward for humans, arithmetic reasoning remains a challenging task for language models \citep[][\textit{inter alia}]{hendrycks2021measuring,patel-etal-2021-nlp}. Remarkably, when applied with the 540B parameter language model, chain-of-thought prompting achieves performance comparable to task-specific fine-tuned models on several tasks, even setting a new state of the art on the difficult GSM8K benchmark \citep{cobbe2021training}.

\subsection{Experimental Setup}

We investigate the use of chain-of-thought prompting across various language models on multiple benchmarks.

\textbf{Benchmarks.}
We evaluate on the following five math word problem datasets:
\textbf{(1)} the \textbf{GSM8K} dataset of math word problems \citep{cobbe2021training}, 
\textbf{(2)} the \textbf{SVAMP} dataset featuring math word problems with diverse structures \citep{patel-etal-2021-nlp},
\textbf{(3)} the \textbf{ASDiv} dataset containing a variety of math word problems \citep{miao-etal-2020-diverse},
\textbf{(4)} the \textbf{AQuA} dataset comprised of algebraic word problems, and
\textbf{(5)} the \textbf{MAWPS} benchmark \citep{koncel-kedziorski-etal-2016-mawps}. Example problems can be found in Appendix \cref{tab:math-datasets}.

\textbf{Standard prompting.} As a baseline, we use standard few-shot prompting as popularized by \citet{brown2020language}, where a language model receives in-context exemplars of input-output pairs before generating an answer for a test example. The exemplars are formatted as questions paired with answers. The model directly outputs the answer, as depicted in \cref{fig:pull-figure} (left).

\textbf{Chain-of-thought prompting.} Our method involves enhancing each example in few-shot prompting by including a chain of thought linked to the corresponding answer, as shown in \cref{fig:pull-figure} (right). Since most datasets only provide an evaluation split, we manually created a collection of eight few-shot exemplars incorporating chains of thought for prompting—the right side of \cref{fig:pull-figure} displays one such chain of thought exemplar, and the complete set is available in Appendix \cref{tab:appendix-mwp-prompts}. (These exemplars were not subject to prompt engineering; the robustness of our approach is analyzed in \cref{subsec:robustness} and \cref{subsec:faq-prompt-engineering}.) To determine if this style of chain-of-thought prompting effectively triggers successful reasoning across various math word problem benchmarks, we applied this single set of eight chain-of-thought exemplars for all benchmarks except AQuA, which uses a multiple-choice format rather than free response. For AQuA, we employed four exemplars with solutions taken from the training set, as listed in Appendix \cref{tab:appendix-aqua-prompt}.

\textbf{Language models.} We assess five large language models. First is \textbf{GPT-3} \citep{brown2020language}, using the versions text-ada-001, text-babbage-001, text-curie-001, and text-davinci-002, which are believed to correspond to InstructGPT models with 350M, 1.3B, 6.7B, and 175B parameters respectively \citep{ouyang2022training}. Second is \textbf{\lamda{}} \citep{thoppilan2022lamda}, available in sizes 422M, 2B, 8B, 68B, and 137B parameters. Third is \textbf{\palm{}}, with models sized at 8B, 62B, and 540B parameters. Fourth is \textbf{UL2 20B} \citep{tay2022unifying}, and fifth is \textbf{Codex} \citep[][code-davinci-002 in the OpenAI API]{chen2021evaluating}. For sampling, we apply greedy decoding (noting that subsequent research indicates chain-of-thought prompting performance can be enhanced by selecting the majority final answer across multiple sampled outputs \citep{wang2022self}). Regarding \lamda{}, results are averaged over five random seeds, each corresponding to a differently randomized order of exemplars. Because \lamda{} showed minimal variance between seeds, to reduce computation, results for all other models are reported using a single exemplar order.

\subsection{Results}
The most notable outcomes of chain-of-thought prompting are summarized in \cref{fig:main-math}, with comprehensive experimental results for each model group, model size, and benchmark provided in \cref{tab:all-lm-math} in the Appendix. Three main conclusions emerge. First, \cref{fig:main-math} demonstrates that chain-of-thought prompting is an ability that arises with increasing model scale \citep{wei2022emergent}. Specifically, chain-of-thought prompting does not enhance performance for smaller models and only provides performance improvements when applied to models with approximately 100B parameters. Our qualitative analysis showed that smaller models generated coherent but illogical chains of thought, resulting in worse performance than standard prompting.

\input{main_fables/main-math}
Second, chain-of-thought prompting yields greater improvements on more complex problems. For example, on GSM8K (which has the lowest baseline performance), the largest GPT and \palm{} models more than doubled their performance. Conversely, for SingleOp, the simplest subset of MAWPS requiring only a single step to solve, performance gains were either negligible or negative (see Appendix \cref{tab:mawps-subsets}).

Third, chain-of-thought prompting using GPT-3 175B and \palm{} 540B competes well with previous state-of-the-art methods, which generally involve finetuning task-specific models on labeled datasets. \cref{fig:main-math} illustrates how \palm{} 540B leverages chain-of-thought prompting to establish new state-of-the-art results on GSM8K, SVAMP, and MAWPS (noting that standard prompting had already surpassed prior best results on SVAMP). For the other two datasets, AQuA and ASDiv, \palm{} with chain-of-thought prompting achieves performance within 2\% of the state of the art (Appendix \cref{tab:all-lm-math}).

To gain a deeper understanding of why chain-of-thought prompting is effective, we conducted a manual review of chains of thought generated by \lamda{} 137B on the GSM8K dataset. Among 50 randomly selected instances where the model produced the correct final answer, every generated chain of thought was logically and mathematically sound except for two cases that, by coincidence, arrived at the correct answer despite errors (refer to \cref{subsec:correct-chain-of-thought} and \cref{tab:appendix-gsm8k-correct-examples} for examples of accurate model-generated chains of thought). Additionally, we analyzed 50 randomly chosen samples where the model outputted incorrect answers. The key findings indicate that 46\% of these chains of thought were nearly correct aside from minor mistakes (such as calculator errors, symbol mapping errors, or a missing reasoning step), while the remaining 54\% contained significant errors related to semantic understanding or coherence (see \cref{subsec:incorrect-chain-of-thought-analysis}). 

To shed light on how scaling enhances chain-of-thought reasoning capabilities, we conducted a comparable error analysis on \palm{} 62B and assessed which errors were resolved by increasing the model size to \palm{} 540B. Our results suggest that scaling \palm{} up to 540B addresses a substantial fraction of one-step omissions and semantic comprehension errors found in the 62B version (see \cref{subsec:why-scale-helps}).

\subsection{Ablation Study}\label{subsec:math-ablation}

The demonstrated advantages of chain-of-thought prompting naturally prompt the question of whether similar performance gains can be achieved through alternative prompting strategies. \cref{fig:bar_ablation} presents an ablation study involving three variations of chain-of-thought prompting detailed below.

\textbf{Equation only.} One hypothesis for the effectiveness of chain-of-thought prompting is that it produces the mathematical equation needed for evaluation. To test this, we experimented with a variant where the model is prompted to generate only the mathematical equation prior to providing the final answer. \cref{fig:bar_ablation} illustrates that equation-only prompting yields limited improvement on GSM8K, suggesting that the semantic complexity of GSM8K questions is too high to be directly converted into an equation without the intermediate natural language reasoning steps characteristic of chain-of-thought prompting. However, for datasets involving one-step or two-step problems, equation-only prompting does enhance performance, as the equation can be straightforwardly extracted from the problem statement (see Appendix \cref{tab:ablations-arithmetic}).

\input{main_fables/ablation-bar}
\textbf{Computation variability alone.} One hypothesis is that chain-of-thought prompting enables the model to allocate more computational effort (i.e., additional intermediate tokens) to challenging problems. To separate the impact of variable computation from that of chain-of-thought reasoning, we experiment with a setup where the model is prompted to output only a sequence of dots ($\ldots$) matching the number of characters in the equation required to solve the problem. This variant achieves performance similar to the baseline, indicating that variable computation alone does not account for the effectiveness of chain-of-thought prompting, and that expressing intermediate steps through natural language seems to provide added value.

\textbf{Chain-of-thought provided after the answer.} Another possible advantage of chain-of-thought prompting is that it might help the model better retrieve relevant knowledge learned during pretraining. To test this, we evaluate a configuration where the chain-of-thought prompt is presented only after the answer, isolating whether the model relies on the generated chain of thought to produce the final answer. This variant performs comparably to the baseline, suggesting that the sequential reasoning reflected in the chain of thought contributes beyond simply activating learned knowledge.

\input{main_fables/main-robustness}
\subsection{Robustness of Chain of Thought}\label{subsec:robustness}

The sensitivity to choice of exemplars is a crucial factor in prompting methods—for example, rearranging the order of few-shot exemplars can cause GPT-3’s accuracy on SST-2 to vary widely, from near random chance (54.3\%) to nearly state-of-the-art levels (93.4\%) \citep{zhao2021calibrate}. In this last subsection, we assess the robustness of chain-of-thought prompting with chains written by different annotators. Besides the earlier results using chains authored by Annotator A, two other co-authors of this paper (Annotators B and C) independently composed chains of thought for the same few-shot examples (displayed in \cref{sec:appendix-alternate-annotators}). Additionally, Annotator A created a more concise version of the chain of thought compared to the original, emulating the style of solutions from \citet{cobbe2021training}.\footnote{For example, while the original chain of thought employs multiple short sentences (\textit{“There were originally 9 computers. For each of 4 days, 5 more computers were added. So 5 * 4 = 20 computers were added. 9 + 20 is 29.”}), the concise chain reads \textit{“5 * 4 = 20 new computers were added. So there are 9 + 20 = 29 new computers in the server room now.”}}

\cref{fig:bar-robustness-analysis} presents these findings for \lamda{} 137B on GSM8K and MAWPS (additional ablation results for other datasets are provided in Appendix \cref{tab:ablations-arithmetic} / \cref{tab:ablations-commonsense-symbolic}). While variability exists across different chain of thought annotations, as anticipated when employing exemplar-based prompting \citep{le-scao-rush-2021-many,reynolds2021prompt,zhao2021calibrate}, every set of chain of thought prompts significantly surpasses the standard baseline. This suggests that effective use of chain of thought does not rely on a specific linguistic style.

To verify that effective chain-of-thought prompting extends to other exemplar sets, we conducted experiments using three sets of eight exemplars randomly drawn from the GSM8K training set, an independent source (the examples in this dataset already contain reasoning steps resembling a chain of thought).\footnote{We selected examples with $\leq 60$ tokens to fit within our input context window, also restricting to $\leq 2$ reasoning steps to fairly compare with the eight exemplars we authored.} \cref{fig:bar-robustness-analysis} indicates that these prompts performed similarly to our manually created exemplars, also greatly outperforming standard prompting.

Beyond robustness to different annotators, independently created chains of thought, diverse exemplars, and various language models, we observe that chain-of-thought prompting for arithmetic reasoning remains stable across different exemplar orders and varying numbers of exemplars (refer to \cref{subsec:faq-prompt-engineering}).

\section{Commonsense Reasoning}\label{sec:commonsense-reasoning}

Although chain of thought is especially effective for math word problems, its language-based framework makes it applicable to a wide range of commonsense reasoning tasks, which require reasoning about physical and human interactions grounded in general background knowledge. Commonsense reasoning is crucial for engaging with the world and remains a challenge for current natural language understanding systems \citep{talmor2022commonsenseqa}.

\textbf{Benchmarks.}  
We utilize five datasets that encompass a wide variety of commonsense reasoning types. The well-known \textbf{CSQA} \citep{talmor-etal-2019-commonsenseqa} presents commonsense questions about the world, involving intricate semantics that frequently depend on prior knowledge.  
\textbf{StrategyQA} \citep{geva-etal-2021-aristotle} challenges models to deduce a multi-step strategy in order to answer questions. From the BIG-bench initiative \citep{bigbench}, we select two specialized evaluation sets: \textbf{Date} Understanding, which entails inferring a date based on given context, and \textbf{Sports} Understanding, which involves judging whether a sports-related sentence is plausible or implausible. Lastly, the \textbf{SayCan} dataset \citep{ahn2022can} focuses on translating a natural language instruction into a sequence of robot actions chosen from a discrete set.  
\cref{fig:dataset-examples} presents examples accompanied by chain of thought annotations for all datasets.  

\textbf{Prompts.}  
We adopt the same experimental procedure as described in the previous section. For CSQA and StrategyQA, we randomly sampled instances from the training set and manually created chains of thought to serve as few-shot exemplars. Since the two BIG-bench tasks lack training sets, we used the first ten examples from the evaluation set as few-shot exemplars and report results on the remaining examples in that set. For SayCan, we employ six examples from the training set used in \citet{ahn2022can}, for which we also manually generated chains of thought.

\textbf{Results.} 
\cref{fig:commonsense-results} presents these findings for \palm{} (complete results for \lamda{}, GPT-3, and various model sizes are provided in \cref{tab:all-lm-commonsense}). 
Across all tasks, increasing the model size enhanced the effectiveness of standard prompting; applying chain-of-thought prompting yielded additional improvements, with the largest gains observed for \palm{} 540B. When using chain-of-thought prompting, \palm{} 540B demonstrated strong performance compared to baselines, surpassing the previous state-of-the-art on StrategyQA (75.6\% versus 69.4\%) and exceeding the performance of an unaided sports enthusiast on sports comprehension (95.4\% versus 84\%). 
These outcomes indicate that chain-of-thought prompting can also boost performance on tasks that require diverse commonsense reasoning capabilities (although improvements were minimal on CSQA).

\input{main_fables/main-commonsense}

\input{main_fables/main-symbolic} 
\section{Symbolic Reasoning}\label{sec:symbolic-reasoning} 
Our final set of experiments addresses symbolic reasoning, which is straightforward for humans but can be difficult for language models. We demonstrate that chain-of-thought prompting not only allows language models to tackle symbolic reasoning tasks that are difficult under standard prompting, but also supports length generalization to inference inputs longer than those included in the few-shot examples.

\paragraph{Tasks.} 
We utilize the following two toy problems.

\begin{itemize}[leftmargin=*,topsep=0pt]
    \itemsep0em \item \textbf{Last letter concatenation.} This task requires the model to join together the last letters of words in a name (for example, \textit{``Amy Brown''} $\rightarrow$ \textit{``yn''}). 
    It is a more difficult variant of first letter concatenation, which language models can already handle without chain of thought.\footnote{We evaluated 10 common names using GPT-3 \texttt{davinci} and it correctly answered all but one.} Full names are generated by randomly combining names from the top one-thousand first and last names obtained from name census data ({\small{\url{https://namecensus.com/}}}).
    \item \textbf{Coin flip.} This task requires the model to determine whether a coin remains heads up after a sequence of people either flip or do not flip the coin (e.g., \textit{``A coin is heads up. Phoebe flips the coin. Osvaldo does not flip the coin. Is the coin still heads up?''} $\rightarrow$ \textit{``no''}).
\end{itemize}

Given the clearly defined nature of these symbolic reasoning tasks, for each task we evaluate on an \textit{in-domain} test set where the examples contain the same number of steps as the training or few-shot exemplars, as well as an \textit{out-of-domain} (OOD) test set where the evaluation examples include more steps than those present in the exemplars. In the case of last letter concatenation, the model is exposed only to exemplars with two-word names and then tested on names comprising 3 and 4 words.\footnote{For names longer than two words, multiple first and last names are concatenated together.} A similar approach is taken for the coin flip task, adjusting the number of potential flips. The experimental design employs the same models and methodologies as described in the preceding sections. Chains of thought for the few-shot exemplars are again manually constructed for each task, as illustrated in \cref{fig:dataset-examples}.

\paragraph{Results.}
\cref{fig:symbolic-main} displays the outcomes of the in-domain and OOD evaluations using \palm{}, while the results for \lamda{} are reported in Appendix \cref{tab:all-lm-symbolic}. With \palm{} 540B, chain-of-thought prompting achieves nearly 100\% accuracy (it is important to note that standard prompting already solves the coin flip task with \palm{} 540B, though this is not the case for \lamda{} 137B). These in-domain evaluations are essentially "toy tasks" because the perfect solution strategies are explicitly provided by the few-shot exemplars’ chains of thought; thus, the model’s role is to replicate the same sequence of steps using new symbols presented in the test examples. Nevertheless, smaller models fail to succeed—competency in performing abstract operations on unseen symbols for these three tasks emerges only when the model size reaches around 100 billion parameters.

Regarding the OOD evaluations, standard prompting fails for both tasks. Using chain-of-thought prompting, language models exhibit upward scaling trends (although their performance remains lower than in the in-domain scenario). Therefore, chain-of-thought prompting enables length generalization beyond previously observed chains of thought for language models of adequate scale.

\section{Discussion}\label{sec:discussion}

We have investigated chain-of-thought prompting as a straightforward technique to evoke multi-step reasoning in large language models. Initially, we observed that chain-of-thought prompting significantly enhances performance on arithmetic reasoning, producing improvements that surpass ablations and remain consistent across different annotators, exemplars, and language models (\cref{sec:arithmetic-reasoning}).
Subsequently, experiments on commonsense reasoning highlighted how the linguistic character of chain-of-thought reasoning renders it broadly applicable (\cref{sec:commonsense-reasoning}).
Lastly, we demonstrated that for symbolic reasoning, chain-of-thought prompting supports OOD generalization to longer sequence lengths (\cref{sec:symbolic-reasoning}).
In all of these experiments, chain-of-thought reasoning is prompted directly from an off-the-shelf language model without any fine-tuning during the course of this study.

The rise of chain-of-thought reasoning as a consequence of increasing model scale has been a central theme \citep{wei2022emergent}. For numerous reasoning tasks where standard prompting shows a flat scaling curve, chain-of-thought prompting produces sharply rising scaling curves. Chain-of-thought prompting appears to broaden the set of tasks that large language models can solve effectively—in other words, our findings emphasize that standard prompting merely establishes a lower bound on the capabilities of large language models. This insight likely prompts more questions than it resolves—for example, how much further improvement in reasoning ability can we anticipate with larger model scales? What other prompting strategies might extend the variety of tasks language models can tackle?

Regarding limitations, we first note that although chain-of-thought mimics human reasoning processes, it does not address whether the neural network is genuinely "reasoning," which remains an open question. Second, while the effort to manually augment exemplars with chains of thought is minimal in a few-shot setting, such annotation costs could be prohibitive for fine-tuning (although this might be mitigated by synthetic data generation or zero-shot generalization).
\orange{Third, there is no assurance that the reasoning paths are correct, which can result in both accurate and inaccurate answers; enhancing the factual reliability of language model generations remains a promising area for future research \citep[][\textit{inter alia}]{rashkin2021measuring,ye2022unreliability,wiegreffe2021reframing}.} Finally, the fact that chain-of-thought reasoning emerges only at large model scales makes it expensive for deployment in practical applications; future work could investigate methods to induce reasoning capabilities in smaller models.

\section{Related Work} This study draws inspiration from several research domains, which we discuss in greater detail in an expanded related work section (\cref{sec:extended-related-work}). Here, we highlight two key directions and their pertinent papers that are particularly relevant.

The first key direction involves utilizing intermediate steps to address reasoning problems. \cite{ling-etal-2017-program} introduce the concept of employing natural language rationales to solve math word problems by breaking down the solution into a sequence of intermediate steps. Their approach stands in notable contrast to prior work that relies on formal languages for reasoning \citep{roy-2016-reasoning, chiang-chen-2019-semantically, amini-etal-2019-mathqa, chen2019neural}. \citet{cobbe2021training} build upon \citet{ling-etal-2017-program} by generating a larger dataset and utilizing it to fine-tune a pretrained language model instead of training a model from scratch. In the area of program synthesis, \cite{nye2021show} employ language models to predict the final outputs of Python programs by first predicting intermediate computational results line-by-line, demonstrating that this stepwise prediction strategy outperforms direct final output prediction.

Unsurprisingly, this paper is also closely connected to the extensive recent research on prompting. Since the advent of few-shot prompting popularized by \citet{brown2020language}, various general techniques have enhanced models' prompting capabilities, such as automatically learning prompts \citep{lester-etal-2021-power} or providing models with task-specific instructions \citep{wei2021finetuned,sanh2021multitask,ouyang2022training}. While these methods improve or modify the input portion of the prompt (e.g., instructions prepended to inputs), our work takes a different approach by augmenting the language models' outputs with a chain of thought.

\section{Conclusions} We have investigated chain-of-thought prompting as a straightforward and widely applicable technique to improve reasoning in language models. Through evaluations on arithmetic, symbolic, and commonsense reasoning tasks, we observe that chain-of-thought reasoning emerges as a property related to model scale, enabling sufficiently large language models to solve reasoning tasks that otherwise exhibit flat scaling performance. Expanding the variety of reasoning tasks that language models can handle may encourage further development of language-based reasoning methods.

\end{document}